% Options for packages loaded elsewhere
\PassOptionsToPackage{unicode}{hyperref}
\PassOptionsToPackage{hyphens}{url}
\documentclass[
]{article}
\usepackage{xcolor}
\usepackage{amsmath,amssymb}
\setcounter{secnumdepth}{-\maxdimen} % remove section numbering
\usepackage{iftex}
\ifPDFTeX
  \usepackage[T1]{fontenc}
  \usepackage[utf8]{inputenc}
  \usepackage{textcomp} % provide euro and other symbols
\else % if luatex or xetex
  \usepackage{unicode-math} % this also loads fontspec
  \defaultfontfeatures{Scale=MatchLowercase}
  \defaultfontfeatures[\rmfamily]{Ligatures=TeX,Scale=1}
\fi
\usepackage{lmodern}
\ifPDFTeX\else
  % xetex/luatex font selection
\fi
% Use upquote if available, for straight quotes in verbatim environments
\IfFileExists{upquote.sty}{\usepackage{upquote}}{}
\IfFileExists{microtype.sty}{% use microtype if available
  \usepackage[]{microtype}
  \UseMicrotypeSet[protrusion]{basicmath} % disable protrusion for tt fonts
}{}
\makeatletter
\@ifundefined{KOMAClassName}{% if non-KOMA class
  \IfFileExists{parskip.sty}{%
    \usepackage{parskip}
  }{% else
    \setlength{\parindent}{0pt}
    \setlength{\parskip}{6pt plus 2pt minus 1pt}}
}{% if KOMA class
  \KOMAoptions{parskip=half}}
\makeatother
\usepackage{color}
\usepackage{fancyvrb}
\newcommand{\VerbBar}{|}
\newcommand{\VERB}{\Verb[commandchars=\\\{\}]}
\DefineVerbatimEnvironment{Highlighting}{Verbatim}{commandchars=\\\{\}}
% Add ',fontsize=\small' for more characters per line
\newenvironment{Shaded}{}{}
\newcommand{\AlertTok}[1]{\textcolor[rgb]{1.00,0.00,0.00}{\textbf{#1}}}
\newcommand{\AnnotationTok}[1]{\textcolor[rgb]{0.38,0.63,0.69}{\textbf{\textit{#1}}}}
\newcommand{\AttributeTok}[1]{\textcolor[rgb]{0.49,0.56,0.16}{#1}}
\newcommand{\BaseNTok}[1]{\textcolor[rgb]{0.25,0.63,0.44}{#1}}
\newcommand{\BuiltInTok}[1]{\textcolor[rgb]{0.00,0.50,0.00}{#1}}
\newcommand{\CharTok}[1]{\textcolor[rgb]{0.25,0.44,0.63}{#1}}
\newcommand{\CommentTok}[1]{\textcolor[rgb]{0.38,0.63,0.69}{\textit{#1}}}
\newcommand{\CommentVarTok}[1]{\textcolor[rgb]{0.38,0.63,0.69}{\textbf{\textit{#1}}}}
\newcommand{\ConstantTok}[1]{\textcolor[rgb]{0.53,0.00,0.00}{#1}}
\newcommand{\ControlFlowTok}[1]{\textcolor[rgb]{0.00,0.44,0.13}{\textbf{#1}}}
\newcommand{\DataTypeTok}[1]{\textcolor[rgb]{0.56,0.13,0.00}{#1}}
\newcommand{\DecValTok}[1]{\textcolor[rgb]{0.25,0.63,0.44}{#1}}
\newcommand{\DocumentationTok}[1]{\textcolor[rgb]{0.73,0.13,0.13}{\textit{#1}}}
\newcommand{\ErrorTok}[1]{\textcolor[rgb]{1.00,0.00,0.00}{\textbf{#1}}}
\newcommand{\ExtensionTok}[1]{#1}
\newcommand{\FloatTok}[1]{\textcolor[rgb]{0.25,0.63,0.44}{#1}}
\newcommand{\FunctionTok}[1]{\textcolor[rgb]{0.02,0.16,0.49}{#1}}
\newcommand{\ImportTok}[1]{\textcolor[rgb]{0.00,0.50,0.00}{\textbf{#1}}}
\newcommand{\InformationTok}[1]{\textcolor[rgb]{0.38,0.63,0.69}{\textbf{\textit{#1}}}}
\newcommand{\KeywordTok}[1]{\textcolor[rgb]{0.00,0.44,0.13}{\textbf{#1}}}
\newcommand{\NormalTok}[1]{#1}
\newcommand{\OperatorTok}[1]{\textcolor[rgb]{0.40,0.40,0.40}{#1}}
\newcommand{\OtherTok}[1]{\textcolor[rgb]{0.00,0.44,0.13}{#1}}
\newcommand{\PreprocessorTok}[1]{\textcolor[rgb]{0.74,0.48,0.00}{#1}}
\newcommand{\RegionMarkerTok}[1]{#1}
\newcommand{\SpecialCharTok}[1]{\textcolor[rgb]{0.25,0.44,0.63}{#1}}
\newcommand{\SpecialStringTok}[1]{\textcolor[rgb]{0.73,0.40,0.53}{#1}}
\newcommand{\StringTok}[1]{\textcolor[rgb]{0.25,0.44,0.63}{#1}}
\newcommand{\VariableTok}[1]{\textcolor[rgb]{0.10,0.09,0.49}{#1}}
\newcommand{\VerbatimStringTok}[1]{\textcolor[rgb]{0.25,0.44,0.63}{#1}}
\newcommand{\WarningTok}[1]{\textcolor[rgb]{0.38,0.63,0.69}{\textbf{\textit{#1}}}}
\usepackage{longtable,booktabs,array}
\newcounter{none} % for unnumbered tables
\usepackage{calc} % for calculating minipage widths
% Correct order of tables after \paragraph or \subparagraph
\usepackage{etoolbox}
\makeatletter
\patchcmd\longtable{\par}{\if@noskipsec\mbox{}\fi\par}{}{}
\makeatother
% Allow footnotes in longtable head/foot
\IfFileExists{footnotehyper.sty}{\usepackage{footnotehyper}}{\usepackage{footnote}}
\makesavenoteenv{longtable}
\setlength{\emergencystretch}{3em} % prevent overfull lines
\providecommand{\tightlist}{%
  \setlength{\itemsep}{0pt}\setlength{\parskip}{0pt}}
\usepackage{bookmark}
\IfFileExists{xurl.sty}{\usepackage{xurl}}{} % add URL line breaks if available
\urlstyle{same}
\hypersetup{
  hidelinks,
  pdfcreator={LaTeX via pandoc}}

\author{}
\date{}

\begin{document}

\section{CIEL: a universal execution engine for distributed data-flow
computing}\label{ciel-a-universal-execution-engine-for-distributed-data-flow-computing}

Derek G. Murray Malte Schwarzkopf Christopher Smowton

Steven Smith Anil Madhavapeddy Steven Hand

University of Cambridge Computer Laboratory

\section{Abstract}\label{abstract}

This paper introduces CIEL, a universal execution engine for distributed
data-flow programs. Like previous execution engines, CIEL masks the
complexity of distributed programming. Unlike those systems, a CIEL job
can make data-dependent control-flow decisions, which enables it to
compute iterative and recursive algorithms.

We have also developed Skywriting, a Turing-complete scripting language
that runs directly on CIEL. The execution engine provides transparent
fault tolerance and distribution to Skywriting scripts and
high-performance code written in other programming languages. We have
deployed CIEL on a cloud computing platform, and demonstrate that it
achieves scalable performance for both iterative and non-iterative
algorithms.

\section{1 Introduction}\label{introduction}

Many organisations have an increasing need to process large data sets,
and a cluster of commodity machines on which to process them.
Distributed execution engines---such as MapReduce {[}18{]} and Dryad
{[}26{]}---have become popular systems for exploiting such clusters.
These systems expose a simple programming model, and automatically
handle the difficult aspects of distributed computing: fault tolerance,
scheduling, synchronisation and communication. MapReduce and Dryad can
be used to implement a wide range of algorithms {[}3, 39{]}, but they
are awkward or inefficient for others {[}12, 21, 25, 28, 34{]}. The
problems typically arise with iterative algorithms, which underlie many
machine-learning and optimisation problems, but require a more
expressive programming model and a more powerful execution engine. To
address these limitations, and extend the benefits of distributed
execution engines to a wider range of applications, we have developed
Skywriting and CIEL.

Skywriting is a scripting language that allows the straightforward
expression of iterative and recursive task-parallel algorithms using
imperative and functional language syntax {[}31{]}. Skywriting scripts
run on CIEL, an execution engine that provides a universal execution
model for distributed data-flow. Like previous systems, CIEL coordinates
the distributed execution of a set of data-parallel tasks arranged
according to a data-flow DAG, and hence benefits from transparent
scaling and fault tolerance. However CIEL extends previous models by
dynamically building the DAG as tasks execute. As we will show, this
conceptually simple extension---allowing tasks to create further
tasks---enables CIEL to support data-dependent iterative or recursive
algorithms. We present the high-level architecture of CIEL in Section 3,
and explain how Skywriting maps onto CIEL's primitives in Section 4.

Our implementation incorporates several additional features, described
in Section 5. Like existing systems, CIEL provides transparent fault
tolerance for worker nodes. Moreover, CIEL can tolerate failures of the
cluster master and the client program. To improve resource utilisation
and reduce execution latency, CIEL can memoise the results of tasks.
Finally, CIEL supports the streaming of data between
concurrently-executing tasks.

We have implemented a variety of applications in Skywriting, including
MapReduce-style (grep, wordcount), iterative (k-means, PageRank) and
dynamic programming (Smith-Waterman, option pricing) algorithms. In
Section 6 we evaluate the performance of some of these applications when
run on a CIEL cluster.

\section{2 Motivation}\label{motivation}

Several researchers have identified limitations in the MapReduce and
Dryad programming models. These systems were originally developed for
batch-oriented jobs, namely large-scale text mining for information
retrieval {[}18, 26{]}. They are designed to maximise throughput, rather
than minimise individual job latency. This is especially noticeable in
iterative computations, for which

Feature

MapReduce{[}2, 18{]}

Dryad{[}26{]}

Pregel{[}28{]}

Iterative MR{[}12, 21{]}

Piccolo{[}34{]}

CIEL

Dynamic control flow

✗

✗

√

√

√

√

Task dependencies

Fixed (2-stage)

Fixed (DAG)

Fixed (BSP)

Fixed (2-stage)

Fixed (1-stage)

Dynamic

Fault tolerance

Transparent

Transparent

Transparent

✗

Checkpoint

Transparent

Data locality

√

√

√

√

√

√

Transparent scaling

√

√

√

√

✗

√

Figure 1: Analysis of the features provided by existing distributed
execution engines.

multiple jobs are chained together and the job latency is multiplied
{[}12, 21, 25, 28, 34{]}.

Nevertheless, MapReduce---in particular its open-source implementation,
Hadoop {[}2{]}---remains a popular platform for parallel iterative
computations with large inputs. For example, the Apache Mahout machine
learning library uses Hadoop as its execution engine {[}3{]}. Several of
the Mahout algorithms---such as k-means clustering and singular value
decomposition---are iterative, comprising a data-parallel kernel inside
a while-not-converged loop. Mahout uses a driver program that submits
multiple jobs to Hadoop and performs convergence testing at the client.
However, since the driver program executes logically (and often
physically) outside the Hadoop cluster, each iteration incurs
job-submission overhead, and the driver program does not benefit from
transparent fault tolerance. These problems are not unique to Hadoop,
but are shared with both the original version of MapReduce {[}18{]} and
Dryad {[}26{]}.

The computational power of a distributed execution engine is determined
by the data flow that it can express. In MapReduce, the data flow is
limited to a bipartite graph parameterised by the number of map and
reduce tasks; Dryad allows data flow to follow a more general directed
acyclic graph (DAG), but it must be fully specified before starting the
job. In general, to support iterative or recursive algorithms within a
single job, we need data-dependent control flow---i.e.~the ability to
create more work dynamically, based on the results of previous
computations. At the same time, we wish to retain the existing benefits
of task-level parallelism: transparent fault tolerance, locality-based
scheduling and transparent scaling. In Figure 1, we analyse a range of
existing systems in terms of these objectives.

MapReduce and Dryad already support transparent fault tolerance,
locality-based scheduling and transparent scaling {[}18, 26{]}. In
addition, Dryad supports arbitrary task dependencies, which enables it
to execute a larger class of computations than MapReduce. However,
neither supports data-dependent control flow, so the work in each
computation must be statically pre-determined.

A variety of systems provide data-dependent control flow but sacrifice
other functionality. Google's Pregel is the largest-scale example of a
distributed execution engine with support for control flow {[}28{]}.
Pregel is a Bulk Synchronous Parallel (BSP) system designed for
executing graph algorithms (such as PageRank), and Pregel computations
are divided into ``supersteps'', during which a ``vertex method'' is
executed for each vertex in the graph. Crucially, each vertex can vote
to terminate the computation, and the computation terminates when all
vertices vote to terminate. Like a simple MapReduce job, however, a
Pregel computation only operates on a single data set, and the
programming model does not support the composition of multiple
computations.

Two recent systems add iteration capabilities to MapReduce.
CGL-MapReduce is a new implementation of MapReduce that caches static
(loop-invariant) data in RAM across several MapReduce jobs {[}21{]}.
HaLoop extends Hadoop with the ability to evaluate a convergence
function on reduce outputs {[}12{]}. Neither system provides fault
tolerance across multiple iterations, and neither can support
Dryad-style task dependency graphs.

Finally, Piccolo is a new programming model for data-parallel
programming that uses a partitioned in-memory key-value table to replace
the reduce phase of MapReduce {[}34{]}. A Piccolo program is divided
into ``kernel'' functions, which are applied to table partitions in
parallel, and typically write key-value pairs into one or more other
tables. A ``control'' function coordinates the kernel functions, and it
may perform arbitrary data-dependent control flow. Piccolo supports
user-assisted checkpointing (based on the Chandy-Lamport algorithm), and
is limited to fixed cluster membership. If a single machine fails, the
entire computation must be restarted from a checkpoint with the same
number of machines.

We believe that CIEL is the first system to support all five goals in
Figure 1, but it is not a panacea. CIEL is designed for coarse-grained
parallelism across large data sets, as are MapReduce and Dryad. For
fine-grained tasks, a work-stealing scheme is more appropriate {[}11{]}.
Where the entire data set can fit in RAM, Piccolo may be more efficient,
because it can avoid writing to disk. Ultimately, achieving the highest
performance requires significant developer effort, using a low-level
technique such as explicit message passing {[}30{]}.

\textit{[Image omitted: images/fe88400b38ed92029c0665ed805053a45576f357e8a770fcd024d0955dd9b18e.jpg]}

flowchart

\begin{Shaded}
\begin{Highlighting}[]
\NormalTok{graph TD}
\NormalTok{    A["u"] {-}{-}\textgreater{} B["A"]}
\NormalTok{    B {-}{-}\textgreater{}|SPAWNS| C["B"]}
\NormalTok{    B {-}{-}\textgreater{}|DELEGATES| D["z"]}
\NormalTok{    C {-}{-}\textgreater{}|x| D}
\NormalTok{    C {-}{-}\textgreater{}|y| D}
\NormalTok{    D {-}{-}\textgreater{} E["Result (future)"]}
\NormalTok{    E {-}{-}\textgreater{} B}
\NormalTok{    F["v"] {-}{-}\textgreater{} C}
\NormalTok{    G["w"] {-}{-}\textgreater{} C}
\NormalTok{    H["Concrete object"] {-}{-}\textgreater{} C}
\NormalTok{    I["Future object"] {-}{-}\textgreater{} D}
\NormalTok{    style A fill:\#f9f,stroke:\#333}
\NormalTok{    style B fill:\#ccf,stroke:\#333}
\NormalTok{    style C fill:\#cfc,stroke:\#333}
\NormalTok{    style D fill:\#fcf,stroke:\#333}
\NormalTok{    style E fill:\#cff,stroke:\#333}
\NormalTok{    style F fill:\#ffc,stroke:\#333}
\NormalTok{    style G fill:\#f9f,stroke:\#333}
\NormalTok{    style H fill:\#ccf,stroke:\#333}
\NormalTok{    style I fill:\#ffc,stroke:\#333}
\end{Highlighting}
\end{Shaded}

\begin{enumerate}
\def\labelenumi{(\alph{enumi})}
\tightlist
\item
  Dynamic task graph
\end{enumerate}

Task ID

Dependencies

Expected outputs

A

\{ u \}

≠

B

\{ v \}

x

C

\{ w \}

y

D

\{ x, y \}

z

Object ID

Produced by

Locations

u

-

\{ host19, host85 \}

v

-

\{ host21, host23 \}

w

-

\{ host22, host57 \}

x

B

∅

y

C

∅

z

A D

∅

\begin{enumerate}
\def\labelenumi{(\alph{enumi})}
\setcounter{enumi}{1}
\tightlist
\item
  Task and object tables\\
  Figure 2: A CIEL job is represented by a dynamic task graph, which
  contains tasks and objects (§3.1). In this example, root task A spawns
  tasks B, C and D, and delegates the production of its result to D.
  Internally, CIEL uses task and object tables to represent the graph
  (§3.3).
\end{enumerate}

\section{3 CIEL}\label{ciel}

CIEL is a distributed execution engine that can execute programs with
arbitrary data-dependent control flow. In this section, we first
describe the core abstraction that CIEL supports: the dynamic task graph
(§3.1). We then describe how CIEL executes a job that is represented as
a dynamic task graph (§3.2). Finally, we describe the concrete
architecture of a CIEL cluster that is used for distributed data-flow
computing (§3.3).

\section{3.1 Dynamic task graphs}\label{dynamic-task-graphs}

In this subsection, we define the three CIEL primitives---objects,
references and tasks---and explain how they are related in a dynamic
task graph (Figure 2).

CIEL is a data-centric execution engine: the goal of a CIEL job is to
produce one or more output objects. An object is an unstructured,
finite-length sequence of bytes. Every object has a unique name: if two
objects exist with the same name, they must have the same contents. To
simplify consistency and replication, an object is immutable once it has
been written, but it is sometimes possible to append to an object
(§5.3).

It is helpful to be able to describe an object without possessing its
full contents; CIEL uses references for this purpose. A reference
comprises a name and a set of locations (e.g.~hostname-port pairs) where
the object with that name is stored. The set of locations may be empty:
in that case, the reference is a future reference to an object that has
not yet been produced. Otherwise, it is a concrete reference, which may
be consumed.

A CIEL job makes progress by executing tasks. A task is a non-blocking
atomic computation that executes completely on a single machine. A task
has one or more dependencies, which are represented by references, and
the task becomes runnable when all of its dependencies become concrete.
The dependencies include a special object that specifies the behaviour
of the task (such as an executable binary or a Java class) and may
impose some structure over the other dependencies. To simplify fault
tolerance (§5.2), CIEL requires that all tasks compute a deterministic
function of their dependencies. A task also has one or more expected
outputs, which are the names of objects that the task will either create
or delegate another task to create.

Tasks can have two externally-observable behaviours. First, a task can
publish one or more objects, by creating a concrete reference for those
objects. In particular, the task can publish objects for its expected
outputs, which may cause other tasks to become runnable if they depend
on those outputs. To support data-dependent control flow, however, a
task may also spawn new tasks that perform additional computation. CIEL
enforces the following conditions on task behaviour:

\begin{enumerate}
\def\labelenumi{\arabic{enumi}.}
\tightlist
\item
  For each of its expected outputs, a task must either publish a
  concrete reference, or spawn a child task with that name as an
  expected output. This ensures that, as long as the children eventually
  terminate, any task that depends on the parent's output will
  eventually become runnable.\\
\item
  A child task must only depend on concrete references (i.e.~objects
  that already exist) or future references to the outputs of tasks that
  have already been spawned (i.e.~objects that are already expected to
  be published). This prevents deadlock, as a cycle cannot form in the
  dependency graph.
\end{enumerate}

The dynamic task graph stores the relation between tasks and objects. An
edge from an object to a task means

that the task depends on that object. An edge from a task to an object
means that the task is expected to output the object. As a job runs, new
tasks are added to the dynamic task graph, and the edges are rewritten
when a newly-spawned task is expected to produce an object.

The dynamic task graph provides low-level data-dependent control flow
that resembles tail recursion: a task either produces its output
(analogous to returning a value) or spawns a new task to produce that
output (analogous to a tail call). It also provides facilities for
data-parallelism, since independent tasks can be dispatched in parallel.
However, we do not expect programmers to construct dynamic task graphs
manually, and instead we provide the Skywriting script language for
generating these graphs programmatically (§4).

\section{3.2 Evaluating objects}\label{evaluating-objects}

Given a dynamic task graph, the role of CIEL is to evaluate one or more
objects that correspond to the job outputs. Indeed, a CIEL job can be
specified as a single root task that has only concrete dependencies, and
an expected output that names the final result of the computation. This
leads to two natural strategies, which are variants of topological
sorting:

Eager evaluation. Since the task dependencies form a DAG, at least one
task must have only concrete dependencies. Start by executing the tasks
with only concrete dependencies; subsequently execute tasks when all of
their dependencies become concrete.

Lazy evaluation. Seek to evaluate the expected output of the root task.
To evaluate an object, identify the task, T, that is expected to produce
the object. If T has only concrete dependencies, execute it immediately;
otherwise, block T and recursively evaluate all of its unfulfilled
dependencies using the same procedure. When the inputs of a blocked task
become concrete, execute it. When the production of a required object is
delegated to a spawned task, re-evaluate that object.

When we first developed CIEL, we experimented with both strategies, but
switched exclusively to lazy evaluation since it more naturally supports
the fault-tolerance and memoisation features that we describe in §5.

\section{3.3 System architecture}\label{system-architecture}

Figure 3 shows the architecture of a CIEL cluster. A single master
coordinates the end-to-end execution of jobs, and several workers
execute individual tasks.

The master maintains the current state of the dynamic task graph in the
object table and task table (Figure 2(b)). Each row in the object table
contains the latest reference for that object, including its locations
(if any), and a pointer to the task that is expected to produce it (if
any: an object will not have a task pointer if it is loaded into the
cluster by an external tool). Each row in the task table corresponds to
a spawned task, and contains pointers to the references on which the
task depends.

\textit{[Image omitted: images/979bb2255af3d6a50f7267a5851b3a765ac5f07fd3357d15d92170bbcde396ef.jpg]}

flowchart

\begin{Shaded}
\begin{Highlighting}[]
\NormalTok{graph LR}
\NormalTok{    A["Master"] {-}{-}\textgreater{}|PUBLISH OBJECT| B["Worker"]}
\NormalTok{    C["Object table"] {-}{-}\textgreater{} D["Scheduler"]}
\NormalTok{    E["Worker table"] {-}{-}\textgreater{} D}
\NormalTok{    F["Task table"] {-}{-}\textgreater{} D}
\NormalTok{    D {-}{-}\textgreater{}|DISPATCH TASK| B}
\NormalTok{    B {-}{-}\textgreater{}|Executors| G["Object store"]}
\NormalTok{    G {-}{-}\textgreater{}|DATA I/O| H["Java"]}
\NormalTok{    G {-}{-}\textgreater{}|DATA I/O| I[".NET"]}
\NormalTok{    G {-}{-}\textgreater{}|DATA I/O| J["..."]}
\NormalTok{    G {-}{-}\textgreater{}|DATA I/O| K["SW"]}
\NormalTok{    D {-}{-}\textgreater{}|SPAWN TASKS| F}
\end{Highlighting}
\end{Shaded}

Figure 3: A CIEL cluster has a single master and many workers. The
master dispatches tasks to the workers for execution. After a task
completes, the worker publishes a set of objects and may spawn further
tasks.

The master scheduler is responsible for making progress in a CIEL
computation: it lazily evaluates output objects and pairs runnable tasks
with idle workers. Since task inputs and outputs may be very large (on
the order of gigabytes per task), all bulk data is stored on the workers
themselves, and the master handles references. The master uses a
multiple-queue-based scheduler (derived from Hadoop {[}2{]}) to dispatch
tasks to the worker nearest the data. If a worker needs to fetch a
remote object, it reads the object directly from another worker.

The workers execute tasks and store objects. At startup, a worker
registers with the master, and periodically sends a heartbeat to
demonstrate its continued availability. When a task is dispatched to a
worker, the appropriate executor is invoked. An executor is a generic
component that prepares input data for consumption and invokes some
computation on it, typically by executing an external process. We have
implemented simple executors for Java, .NET, shell-based and native
code, as well as a more complex executor for Skywriting (§4).

Assuming that a worker executes a task successfully, it will reply to
the master with the set of references that it wishes to publish, and a
list of task descriptors for any new tasks that it wishes to spawn. The
master will then update the object table and task table, and re-evaluate
the set of tasks now runnable.

In addition to the master and workers, there will be one or more clients
(not shown). A client's role is minimal: it submits a job to the master,
and either polls the master to discover the job status or blocks until
the job completes.

\begin{Shaded}
\begin{Highlighting}[]
\KeywordTok{function} \VariableTok{process\_chunk}\NormalTok{(}\VariableTok{chunk}\OperatorTok{,} \VariableTok{prev\_result}\NormalTok{) \{}
    \OperatorTok{//} \VariableTok{Execute} \VariableTok{native} \VariableTok{code} \KeywordTok{for} \VariableTok{chunk} \VariableTok{processing}\NormalTok{.}
    \OperatorTok{//} \VariableTok{Returns} \VariableTok{a} \VariableTok{reference} \VariableTok{to} \VariableTok{a} \VariableTok{partial} \VariableTok{result}\NormalTok{.}
    \KeywordTok{return} \VariableTok{spawn\_exec}\NormalTok{(}\OperatorTok{...}\NormalTok{)}\OperatorTok{;}
\NormalTok{\}}

\KeywordTok{function} \VariableTok{is\_converged}\NormalTok{(}\VariableTok{curr\_result}\OperatorTok{,} \VariableTok{prev\_result}\NormalTok{) \{}
    \OperatorTok{//} \VariableTok{Execute} \VariableTok{native} \VariableTok{code} \KeywordTok{for} \VariableTok{convergence} \VariableTok{test}\NormalTok{.}
    \OperatorTok{//} \VariableTok{Returns} \VariableTok{a} \VariableTok{reference} \VariableTok{to} \VariableTok{a} \VariableTok{boolean}\NormalTok{.}
    \KeywordTok{return} \VariableTok{spawn\_exec}\NormalTok{(}\OperatorTok{...}\NormalTok{)[}\FloatTok{0}\NormalTok{]}\OperatorTok{;}
\NormalTok{\}}

\VariableTok{input\_data} \OperatorTok{=}\NormalTok{ [}\VariableTok{ref}\NormalTok{(}\StringTok{"ciel://host137/chunk0"}\NormalTok{)}\OperatorTok{,}
    \VariableTok{ref}\NormalTok{(}\StringTok{"ciel://host223/chunk1"}\NormalTok{)}\OperatorTok{,}
    \OperatorTok{...}\NormalTok{]}\OperatorTok{;}
\VariableTok{curr} \OperatorTok{=} \OperatorTok{...;} \OperatorTok{//} \VariableTok{Initial} \VariableTok{guess} \VariableTok{at} \VariableTok{the} \VariableTok{result}\NormalTok{.}

\VariableTok{do}\NormalTok{ \{}
    \VariableTok{prev} \OperatorTok{=} \VariableTok{curr}\OperatorTok{;}
    \VariableTok{curr} \OperatorTok{=}\NormalTok{ []}\OperatorTok{;}
    \KeywordTok{for}\NormalTok{ (}\VariableTok{chunk} \VariableTok{in} \VariableTok{input\_data}\NormalTok{) \{}
    \VariableTok{curr} \OperatorTok{+=} \VariableTok{process\_chunk}\NormalTok{(}\VariableTok{chunk}\OperatorTok{,} \VariableTok{prev}\NormalTok{)}\OperatorTok{;}
\NormalTok{    \}}
\NormalTok{\} }\KeywordTok{while}\NormalTok{ (}\BaseNTok{!*is\_converged(curr, prev));}

\KeywordTok{return} \VariableTok{curr}\OperatorTok{;} 
\end{Highlighting}
\end{Shaded}

Figure 4: Iterative computation implemented in Skywriting. input\_data
is a list of n input chunks, and curr is initialised to a list of n
partial results.

A job submission message contains a root task, which must have only
concrete dependencies. The master adds the root task to the task table,
and starts the job by lazily evaluating its output (§3.2).

Note that CIEL currently uses a single (active) master for simplicity.
Despite this, our implementation can recover from master failure (
\(\S5.2\) ), and it did not cause a performance bottleneck during our
evaluation ( \(\S6\) ). Nonetheless, if it became a concern in future,
it would be possible to partition the master state---i.e.~the task table
and object table---between several hosts, while retaining the
functionality of a single logical master.

\section{4 Skywriting}\label{skywriting}

Skywriting is a language for expressing task-level parallelism that runs
on top of CIEL. Skywriting is Turing-complete, and can express arbitrary
data-dependent control flow using constructs such as while loops and
recursive functions. Figure 4 shows an example Skywriting script that
computes an iterative algorithm; we use a similar structure in the
k-means experiment (§6.2).

We introduced Skywriting in a previous paper {[}31{]}, but briefly
restate the key features here:

- ref(url) returns a reference to the data stored at the given URL. The
function supports common URL schemes, and the custom ciel scheme, which
accesses entries in the CIEL object table. If the URL is external, CIEL
downloads the data into the cluster as an object, and assigns a name for
the object.

\textit{[Image omitted: images/66699d50e2b562b9bdfd1881ce4e8a4b6ae7ad9a0c31aca7108519a1a2760414.jpg]}

flowchart

\begin{Shaded}
\begin{Highlighting}[]
\NormalTok{graph TD}
\NormalTok{    A["function f(x) \{ ... \} return f(42);"] {-}{-}\textgreater{} B["T"]}
\NormalTok{    B {-}{-}\textgreater{} C["t₀"] {-}{-}\textgreater{} D["result"]}
\end{Highlighting}
\end{Shaded}

\begin{enumerate}
\def\labelenumi{(\alph{enumi})}
\tightlist
\item
  Skywriting task
\end{enumerate}

\textit{[Image omitted: images/2f7dc131570e979942ed7a2f1ff845857aed97fd001a53b89fc37af373eb8a0b.jpg]}

flowchart

\begin{Shaded}
\begin{Highlighting}[]
\NormalTok{graph TD}
\NormalTok{    A["jar = z\textbackslash{}ninputs = x, y\textbackslash{}ncls = a.b.Foo"] {-}{-}\textgreater{} C["T"]}
\NormalTok{    B["x"] {-}{-}\textgreater{} C}
\NormalTok{    D["y"] {-}{-}\textgreater{} C}
\NormalTok{    E["z"] {-}{-}\textgreater{} C}
\NormalTok{    C {-}{-}\textgreater{} F["t₁"]}
\NormalTok{    C {-}{-}\textgreater{} G["..."]}
\NormalTok{    C {-}{-}\textgreater{} H["tₙ"]}
\NormalTok{    style C fill:\#f9f,stroke:\#333}
\NormalTok{    style F stroke{-}dasharray: 5 5}
\NormalTok{    style G stroke{-}dasharray: 5 5}
\NormalTok{    style H stroke{-}dasharray: 5 5}
\NormalTok{    note right of C}
\NormalTok{        n results}
\NormalTok{    end}
\end{Highlighting}
\end{Shaded}

\begin{enumerate}
\def\labelenumi{(\alph{enumi})}
\setcounter{enumi}{1}
\tightlist
\item
  Other (e.g.~Java) tasks
\end{enumerate}

\textit{[Image omitted: images/03defd1f15096ad3d9946137b5462d06910766432cedae6bf95408e087cd23a5.jpg]}

flowchart

\begin{Shaded}
\begin{Highlighting}[]
\NormalTok{graph TD}
\NormalTok{    A["a = spawn(f); b = spawn(g); return *a + *b"] {-}{-}\textgreater{} B["T"]}
\NormalTok{    B {-}{-}\textgreater{} C["t₀"]}
\NormalTok{    D["f(); g();"] {-}{-}\textgreater{} E["F"]}
\NormalTok{    D {-}{-}\textgreater{} F["G"]}
\NormalTok{    E {-}{-}\textgreater{} G["T\textquotesingle{}"]}
\NormalTok{    F {-}{-}\textgreater{} G}
\NormalTok{    G {-}{-}\textgreater{} H["Continuation of T"]}
\NormalTok{    H {-}{-}\textgreater{} I["a = spawn(f); b = spawn(g); return *a + *b"]}
\end{Highlighting}
\end{Shaded}

\begin{enumerate}
\def\labelenumi{(\alph{enumi})}
\setcounter{enumi}{2}
\tightlist
\item
  Implicit continuation due to dereferencing\\
  Figure 5: Task creation in Skywriting. Tasks can be created using (a)
  spawn(), (b) spawn\_exec() and (c) the dereference (*) operator.
\end{enumerate}

\begin{itemize}
\tightlist
\item
  spawn(f, {[}arg, \ldots{]}) spawns a parallel task to evaluate f(arg,
  \ldots). Skywriting functions are pure: functions cannot have
  side-effects, and all arguments are passed by value. The return value
  is a reference to the result of f(arg, \ldots).\\
\item
  exec(executor, args, n) synchronously runs the named executor with the
  given args. The executor will produce n outputs. The return value is a
  list of n references to those outputs.\\
\item
  spawn\_exec(executor, args, n) spawns a parallel task to run the named
  executor with the given args. As with exec(), the return value is a
  list of n references to those outputs.\\
\item
  The dereference (unary- \(^{*}\) ) operator can be applied to any
  reference; it loads the referenced data into the Skywriting execution
  context, and evaluates to the resulting data structure.
\end{itemize}

In the following, we describe how Skywriting maps on to CIEL primitives.
We describe how tasks are created (§4.1), how references are used to
facilitate data-dependent control flow (§4.2), and the relationship
between Skywriting and other frameworks (§4.3).

\section{4.1 Creating tasks}\label{creating-tasks}

The distinctive feature of Skywriting is its ability to spawn new tasks
in the middle of executing a job. The language provides two explicit
mechanisms for spawning new tasks (the spawn() and spawn\_exec()
functions) and one implicit mechanism (the \(\star\) -operator). Figure
5 summarises these mechanisms.

The spawn() function creates a new task to run the given Skywriting
function. To do this, the Skywriting runtime first creates a data object
that contains the new task's environment, including the text of the
function to be executed and the values of any arguments passed to the
function. This object is called a Skywriting continuation, because it
encapsulates the state of a computation. The runtime then creates a task
descriptor for the new task, which includes a dependency on the new
continuation. Finally, it assigns a reference for the task result, which
it returns to the calling script. Figure 5(a) shows the structure of the
created task.

The spawn\_exec() function is a lower-level task-creation mechanism that
allows the caller to invoke code written in a different language.
Typically, this function is not called directly, but rather through a
wrapper for the relevant executor (e.g.~the built-in java() library
function). When spawn\_exec() is called, the runtime serialises the
arguments into a data object and creates a task that depends on that
object (Figure 5(b)). If the arguments to spawn\_exec() include
references, the runtime adds those references to the new task's
dependencies, to ensure that CIEL will not schedule the task until all
of its arguments are available. Again, the runtime creates references
for the task outputs, and returns them to the calling script. We discuss
how names are chosen in §5.1.

If the task attempts to dereference an object that has not yet been
created---for example, the result of a call to spawn()---the current
task must block. However, CIEL tasks are non-blocking: all
synchronisation (and data-flow) must be made explicit in the dynamic
task graph (§3.1). To resolve this contradiction, the runtime implicitly
creates a continuation task that depends on the dereferenced object and
the current continuation (i.e.~the current Skywriting execution stack).
The new task therefore will only run when the dereferenced object has
been produced, which provides the necessary synchronisation. Figure 5(c)
shows the dependency graph that results when a task dereferences the
result of spawn().

A task terminates when it reaches a return statement (or it blocks on a
future reference). A Skywriting task has a single output, which is the
value of the expression in the return statement. On termination, the
runtime stores the output in the local object store, publishes a
concrete reference to the object, and sends a list of spawned tasks to
the master, in order of creation. Skywriting ensures that the dynamic
task graph remains acyclic. A task's dependencies are fixed when the
task-creation function is evaluated, which means that they can only
include references that are stored in the local Skywriting scope before
evaluating the function. Therefore, a task cannot depend on itself or
any of its descendants. Note that the results of spawn() and
spawn\_exec() are first-class futures {[}24{]}: a Skywriting task can
pass the references in its return value or in a subsequent call to the
task-creation functions. This enables a script to create arbitrary
acyclic dependency graphs, such as the MapReduce dependency graph
(§4.3).

\section{4.2 Data-dependent control
flow}\label{data-dependent-control-flow}

Skywriting is designed to coordinate data-centric computations, which
means that the objects in the computation can be divided into two
spaces:

Data space. Contains large data objects that may be up to several
gigabytes in size.

Coordination space. Contains small objects---such as integers, booleans,
strings, lists and dictionaries---that determine the control flow.

In general, objects in the data space are processed by programs written
in compiled languages, to achieve better I/O or computational
performance than Skywriting can provide. In existing distributed
execution engines (such as MapReduce and Dryad), the data space and
coordination space are disjoint, which prevents these systems from
supporting data-dependent control flow.

To support data-dependent control flow, data must be able to pass from
the data space into the coordination space, so that it can help to
determine the control flow. In Skywriting, the *-operator transforms a
reference to a (data space) object into a (coordination space) value.
The producing task, which may be run by any executor, must write the
referenced object in a format that Skywriting can recognise; we use
JavaScript Object Notation (JSON) for this purpose {[}4{]}. This
serialisation format is only used for references that are passed to
Skywriting, and the majority of executors use the appropriate binary
format for their data.

\section{4.3 Other languages and
frameworks}\label{other-languages-and-frameworks}

Systems like MapReduce have become popular, at least in part, because of
their simple interface: a developer can specify a whole distributed
computation with just a pair of map() and reduce() functions. To
demonstrate that Skywriting approaches this level of simplicity, Figure
6 shows an implementation of the MapReduce execution model, taken from
the Skywriting standard library.

\begin{Shaded}
\begin{Highlighting}[]
\KeywordTok{function} \FunctionTok{apply}\NormalTok{(f}\OperatorTok{,}\NormalTok{ list) \{}
\NormalTok{    outputs }\OperatorTok{=}\NormalTok{ []}\OperatorTok{;}
    \ControlFlowTok{for}\NormalTok{ (i }\KeywordTok{in} \FunctionTok{range}\NormalTok{(}\FunctionTok{len}\NormalTok{(list))) \{}
\NormalTok{    outputs[i] }\OperatorTok{=} \FunctionTok{f}\NormalTok{(list[i])}\OperatorTok{;}
\NormalTok{    \}}
    \ControlFlowTok{return}\NormalTok{ outputs}\OperatorTok{;}
\NormalTok{\}}

\KeywordTok{function} \FunctionTok{shuffle}\NormalTok{(inputs}\OperatorTok{,}\NormalTok{ num\_outputs) \{}
\NormalTok{    outputs }\OperatorTok{=}\NormalTok{ []}\OperatorTok{;}
    \ControlFlowTok{for}\NormalTok{ (i }\KeywordTok{in} \FunctionTok{range}\NormalTok{(num\_outputs)) \{}
\NormalTok{    outputs[i] }\OperatorTok{=}\NormalTok{ []}\OperatorTok{;}
    \ControlFlowTok{for}\NormalTok{ (j }\KeywordTok{in} \FunctionTok{range}\NormalTok{(}\FunctionTok{len}\NormalTok{(inputs))) \{}
\NormalTok{    outputs[i][j] }\OperatorTok{=}\NormalTok{ inputs[j][i]}\OperatorTok{;}
\NormalTok{    \}}
\NormalTok{    \}}
    \ControlFlowTok{return}\NormalTok{ outputs}\OperatorTok{;}
\NormalTok{\}}

\KeywordTok{function} \FunctionTok{mapreduce}\NormalTok{(inputs}\OperatorTok{,}\NormalTok{ mapper}\OperatorTok{,}\NormalTok{ reducer}\OperatorTok{,}\NormalTok{ r) \{}
\NormalTok{    map\_outputs }\OperatorTok{=} \FunctionTok{apply}\NormalTok{(mapper}\OperatorTok{,}\NormalTok{ inputs)}\OperatorTok{;}
\NormalTok{    reduce\_inputs }\OperatorTok{=} \FunctionTok{shuffle}\NormalTok{(map\_outputs}\OperatorTok{,}\NormalTok{ r)}\OperatorTok{;}
\NormalTok{    reduce\_outputs }\OperatorTok{=} \FunctionTok{apply}\NormalTok{(reducer}\OperatorTok{,}\NormalTok{ reduce\_inputs)}\OperatorTok{;}
    \ControlFlowTok{return}\NormalTok{ reduce\_outputs}\OperatorTok{;}
\NormalTok{\} }
\end{Highlighting}
\end{Shaded}

Figure 6: Implementation of the MapReduce programming model in
Skywriting. The user provides a list of inputs, a mapper function, a
reducer function and the number of reducers to use.

The mapreduce() function first applies mapper to each element of inputs.
mapper is a Skywriting function that returns a list of r elements. The
map outputs are then shuffled, so that the \(i^{\text{th}}\) output of
each map becomes an input to the \(i^{\text{th}}\) reduce. Finally, the
reducer function is applied r times to the collected reduce inputs. In
typical use, the inputs to mapreduce() are data objects containing the
input splits, and the mapper and reducer functions invoke spawn\_exec()
to perform computation in another language.

Note that the mapper function is responsible for partitioning data
amongst the reducers, and the reducer function must merge the inputs
that it receives. The implementation of mapper may also incorporate a
combiner, if desired {[}18{]}. To simplify development, we have ported
portions of the Hadoop MapReduce framework to run as CIEL tasks, and
provide helper functions for partitioning, merging, and processing
Hadoop file formats.

Any higher-level language that is compiled into a DAG of tasks can also
be compiled into a Skywriting program, and executed on a CIEL cluster.
For example, one could develop Skywriting back-ends for Pig {[}32{]} and
DryadLINQ {[}39{]}, raising the possibility of extending those languages
with support for unbounded iteration.

\section{5 Implementation issues}\label{implementation-issues}

The current implementation of CIEL and Skywriting contains approximately
9,500 lines of Python code, and a few hundred lines of C, Java and other
languages in the executor bindings. All of the source code, along with a
suite of example Skywriting programs (including those used to evaluate
the system in §6), is available to download from our project website:

http://www.cl.cam.ac.uk/netos/ciel/

The remainder of this section describes three interesting features of
our implementation: memoisation (§5.1), master fault tolerance (§5.2)
and streaming (§5.3).

\section{5.1 Deterministic naming \&
memoisation}\label{deterministic-naming-memoisation}

Recall that all objects in a CIEL cluster have a unique name. In this
subsection, we show how an appropriate choice of names can enable
memoisation.

Our original implementation of CIEL used globally-unique identifiers
(UUIDs) to identify all data objects. While this was a conceptually
simple scheme, it complicated fault tolerance (see following
subsection), because the master had to record the generated UUIDs to
support deterministic task replay after a failure.

This motivated us to reconsider the choice of names. To support
fault-tolerance, existing systems assume that individual tasks are
deterministic {[}18, 26{]}, and CIEL makes the same assumption (§3.1).
It follows that two tasks with the same dependencies---including the
executable code as a dependency---will have identical behaviour.
Therefore the n outputs of a task created with the following Skywriting
statement

\begin{Shaded}
\begin{Highlighting}[]
\NormalTok{result = spawn\_exec(executor, args, n); }
\end{Highlighting}
\end{Shaded}

will be completely determined by executor, args, n and their indices. We
could therefore construct a name for the \(i^{th}\) output by
concatenating executor, args, n and i, with appropriate delimiters.
However, since args may itself contain references, names could grow to
an unmanageable length. We therefore use a collision-resistant hash
function, H, to compute a digest of args and n, which gives the
resulting name:

\begin{Shaded}
\begin{Highlighting}[]
\NormalTok{executor : | H(args||n) : | i }
\end{Highlighting}
\end{Shaded}

We currently use the 160-bit SHA-1 hash function to generate the digest.

Recall the lazy evaluation algorithm from §3.2: tasks are only executed
when their expected outputs are needed to resolve a dependency for a
blocked task. If a new task's outputs have already been produced by a
previous task, the new task need not be executed at all. Hence, as a
result of deterministic naming, CIEL memoises task results, which can
improve the performance of jobs that perform repetitive tasks.

The goals of our memoisation scheme are similar to the recent Nectar
system {[}23{]}. Nectar performs static

analysis on DryadLINQ queries to identify subqueries that have
previously been computed on the same data. Nectar is implemented at the
DryadLINQ level, which enables it to make assumptions about the
semantics of the each task, and the cost/benefit ratio of caching
intermediate results. For example, Nectar can re-use the results of
commutative and associative aggregations from a previous query, if the
previous query operated on a prefix of the current query's input. The
expressiveness of CIEL jobs makes it more challenging to run these
analyses, and we are investigating how simple annotations in a
Skywriting program could provide similar functionality in our system.

\section{5.2 Fault tolerance}\label{fault-tolerance}

A distributed execution engine must continue to make progress in the
face of network and computer faults. As jobs become longer---and, since
CIEL allows unbounded iteration, they may become extremely long---the
probability of experiencing a fault increases. Therefore, CIEL must
tolerate the failure of any machine involved in the computation: the
client, workers and master.

Client fault tolerance is trivial, since CIEL natively supports
iterative jobs and manages job execution from start to finish. The
client's only role is to submit the job: if the client subsequently
fails, the job will continue without interruption. By contrast, in order
to execute an iterative job using a non-iterative framework, the client
must run a driver program that performs all data-dependent control flow
(such as convergence testing). Since the driver program executes outside
the framework, it does not benefit from transparent fault tolerance, and
the developer must provide this manually, for example by checkpointing
the execution state. In our system, a Skywriting script replaces the
driver program, and CIEL executes the whole script reliably.

Worker fault tolerance in CIEL is similar to Dryad {[}26{]}. The master
receives periodic heartbeat messages from each worker, and considers a
worker to have failed if (i) it has not sent a heartbeat after a
specified timeout, and (ii) it does not respond to a reverse message
from the master. At this point, if the worker has been assigned a task,
that task is deemed to have failed.

When a task fails, CIEL automatically re-executes it. However, if it has
failed because its inputs were stored on a failed worker, the task is no
longer runnable. In that case, CIEL recursively re-executes predecessor
tasks until all of the failed task's dependencies are resolved. To
achieve this, the master invalidates the locations in the object table
for each missing input, and lazily re-evaluates the missing inputs.
Other tasks that depend on data from the failed worker will also fail,
and these are similarly re-executed by the master.

Master fault tolerance is also supported in CIEL. In MapReduce and
Dryad, a job fails completely if its master process fails {[}18, 26{]};
in Hadoop, all jobs fail if the JobTracker fails {[}2{]}; and master
failure will usually cause driver programs that submit multiple jobs to
fail. However, in CIEL, all master state can be derived from the set of
active jobs. At a minimum, persistently storing the root task of each
active job allows a new master to be created and resume execution
immediately. CIEL provides three complementary mechanisms that extend
master fault tolerance: persistent logging, secondary masters and object
table reconstruction.

When a new job is created, the master creates a log file for the job,
and synchronously writes its root task descriptor to the log. By
default, it writes the log to a log directory on local secondary
storage, but it can also write to a networked file system or distributed
storage service. As new tasks are created, their descriptors are
appended asynchronously to the log file, and periodically flushed to
disk. When the job completes, a concrete reference to its result is
written to the log directory. Upon restarting, the master scans its log
directory for jobs without a matching result. For those jobs, it replays
the log, rebuilding the dynamic task graph, and ignoring the final
record if it is truncated. Once all logs have been processed, the master
restarts the jobs by lazily evaluating their outputs.

Alternatively, the master may log state updates to a secondary master.
After the secondary master registers with the primary master, the
primary asynchronously forwards all task table and object table updates
to the secondary. Each new job is sent synchronously, to ensure that it
is logged at the secondary before the client receives an
acknowledgement. In addition, the secondary records the address of every
worker that registers with the primary, so that it can contact the
workers in a fail-over scenario. The secondary periodically sends a
heartbeat to the primary; when it detects that the primary has failed,
the secondary instructs all workers to re-register with it. We evaluate
this scenario in §6.5.

If the master fails and subsequently restarts, the workers can help to
reconstruct the object table using the contents of their local object
stores. A worker deems the master to have failed if it does not respond
to requests. At this point, the worker switches into reregister mode,
and the heartbeat messages are replaced with periodic registration
requests to the same network location. When the worker finally contacts
a new master, the master pulls a list of the worker's data objects,
using a protocol based on GFS master recovery {[}22{]}.

\section{5.3 Streaming}\label{streaming}

Our earlier definition of a task (§3.1) stated that a task produces data
objects as part of its result. This definition

implies that object production is atomic: an object either exists
completely or not at all. However, since data objects may be very large,
there is often the opportunity to stream the partially-written object
between tasks, which can lead to pipelined parallelism.

If the producing task has streamable outputs, it sends a pre-publish
message to the master, containing stream references for each streamable
output. These references are used to update the object table, and may
unblock other tasks: the stream consumers. A stream consumer executes as
before, but the executed code reads its input from a named pipe rather
than a local file. A separate thread in the consuming worker process
fetches chunks of input from the producing worker, and writes them into
the pipe. When the producer terminates successfully, it commits its
outputs, which signals to the consumer that no more data remains to be
read.

In the present implementation, the stream producer also writes its
output data to a local disk, so that, if the stream consumer fails, the
producer is unaffected. If the producer fails while it has a consumer,
the producer rolls back any partially-written output. In this case, the
consumer will fail due to missing input, and trigger re-execution of the
producer (§5.2). We are investigating more sophisticated fault-tolerance
and scheduling policies that would allow the producer and consumer to
communicate via direct TCP streams, as in Dryad {[}26{]} and the Hadoop
Online Prototype {[}16{]}. However, as we show in the following section,
support for streaming yields useful performance benefits for some
applications.

\section{6 Evaluation}\label{evaluation}

Our main goal in developing CIEL was to develop a system that supports a
more powerful model of computation than existing distributed execution
engines, without incurring a high cost in terms of performance. In this
section, we evaluate the performance of CIEL running a variety of
applications implemented in Skywriting. We investigate the following
questions:

\begin{enumerate}
\def\labelenumi{\arabic{enumi}.}
\tightlist
\item
  How does CIEL's performance compare to a system in production use
  (viz.~Hadoop)? (§6.1, §6.2)\\
\item
  What benefits does CIEL provide when executing an iterative algorithm?
  (§6.2)\\
\item
  What overheads does CIEL impose on compute-intensive tasks? (§6.3,
  §6.4)\\
\item
  What effect does master failure have on end-to-end job performance?
  (§6.5)
\end{enumerate}

For our evaluation, we selected a set of algorithms to answer these
questions, including MapReduce-style, iterative, and compute-intensive
algorithms. We chose dynamic programming algorithms to demonstrate
CIEL's ability to execute algorithms with data dependencies that do not
translate to the MapReduce model.

\textit{[Image omitted: images/4e3b2b5466ea995136045581e7a52448b2d0cb0916341b6e54d4df55092f4e50.jpg]}

bar

{\def\LTcaptype{none} % do not increment counter
\begin{longtable}[]{@{}lll@{}}
\toprule\noalign{}
Number of workers & Hadoop (s) & CIEL (s) \\
\midrule\noalign{}
\endhead
\bottomrule\noalign{}
\endlastfoot
10 & 350 & 250 \\
20 & 220 & 145 \\
50 & 135 & 85 \\
100 & 90 & 55 \\
\end{longtable}
}

Figure 7: Grp execution time on Hadoop and CIEL (§6.1).

All of the results presented in this section were gathered using
m1.small virtual machines on the Amazon EC2 cloud computing platform. At
the time of writing, an m1.small instance has 1.7 GB of RAM and 1
virtual core (equivalent to a 2007 AMD Opteron or Intel Xeon processor)
{[}1{]}. In all cases, the operating system was Ubuntu 10.04, using
Linux kernel version 2.6.32 in 32-bit mode. Since the virtual machines
are single-core, we run one CIEL worker per machine, and configure
Hadoop to use one map slot per TaskTracker.

\section{6.1 Grep}\label{grep}

Our grep benchmark uses the Grep example application from Hadoop to
search a 22.1 GB dump of English-language Wikipedia for a
three-character string. The original Grep application performs two
MapReduce jobs: the first job parses the input data and emits the
matching strings, and the second sorts the matching strings by
frequency. In Skywriting, we implemented this as a single script that
uses two invocations of mapreduce() (§4.3). Both systems use identical
data formats and execute an identical computation (regular expression
matching).

Figure 7 shows the absolute execution time for Grep as the number of
workers increases from 10 to 100. Averaged across all runs, CIEL
outperforms Hadoop by 35\%. We attribute this to the Hadoop heartbeat
protocol, which limits the rate at which TaskTrackers poll for tasks
once every 5 seconds, and the mandatory ``setup'' and ``cleanup'' phases
that run at the start and end of each job {[}38{]}. As a result, the
relative performance of CIEL improves as the job becomes shorter: CIEL
takes 29\% less time on 10 workers, and 40\% less time on 100

\textit{[Image omitted: images/998371d88097360786deb985ffe8d692dd14ee640417af5a009c36192fdc1f40.jpg]}

line

{\def\LTcaptype{none} % do not increment counter
\begin{longtable}[]{@{}lll@{}}
\toprule\noalign{}
Number of tasks & Hadoop & CIEL \\
\midrule\noalign{}
\endhead
\bottomrule\noalign{}
\endlastfoot
20 & 250 & 150 \\
40 & 450 & 300 \\
60 & 600 & 450 \\
80 & 750 & 600 \\
100 & 900 & 750 \\
\end{longtable}
}

\begin{enumerate}
\def\labelenumi{(\alph{enumi})}
\tightlist
\item
  Iteration length
\end{enumerate}

\textit{[Image omitted: images/8c341aa0eaa776bd988220752abcb76bda10715a46720106709d5b5838c29955.jpg]}

line

{\def\LTcaptype{none} % do not increment counter
\begin{longtable}[]{@{}lll@{}}
\toprule\noalign{}
Time since start (s) & Hadoop & CIEL \\
\midrule\noalign{}
\endhead
\bottomrule\noalign{}
\endlastfoot
0 & 100\% & 100\% \\
3751 & 0\% & 0\% \\
4539 & 100\% & 100\% \\
\end{longtable}
}

\begin{enumerate}
\def\labelenumi{(\alph{enumi})}
\setcounter{enumi}{1}
\tightlist
\item
  Cluster utilisation (100 tasks)
\end{enumerate}

\textit{[Image omitted: images/8b4ecd79df1e37407d583d6f304fdb7dd9b862d6a4b57d6defaed8509f3fe5ff.jpg]}

line

{\def\LTcaptype{none} % do not increment counter
\begin{longtable}[]{@{}lll@{}}
\toprule\noalign{}
Task duration (s) & CIEL & Hadoop \\
\midrule\noalign{}
\endhead
\bottomrule\noalign{}
\endlastfoot
0 & 0 & 0 \\
50 & 0 & 0 \\
100 & 0 & 0 \\
150 & 1 & 0.5 \\
200 & 1 & 1 \\
\end{longtable}
}

\begin{enumerate}
\def\labelenumi{(\alph{enumi})}
\setcounter{enumi}{2}
\tightlist
\item
  Map task distribution\\
  Figure 8: Results of the k-means experiment on Hadoop and CIEL with 20
  workers (§6.2).
\end{enumerate}

workers. We observed that a no-op Hadoop job (which dispatches one map
task per worker, and terminates immediately) runs for an average of 30
seconds. Since Grep involves two jobs, we would not expect Hadoop to
complete the benchmark in less than 60 seconds. These results confirm
that Hadoop is not well-suited to short jobs, which is a result of its
original application (large-scale document indexing). However, anecdotal
evidence suggests that production Hadoop clusters mostly run jobs
lasting less than 90 seconds {[}40{]}.

\section{6.2 k-means}\label{k-means}

We ported the Hadoop-based k-means implementation from the Apache Mahout
scalable machine learning toolkit {[}3{]} to CIEL. Mahout simulates
iterative-algorithm support on Hadoop by submitting a series of jobs and
performing a convergence test outside the cluster; our port uses a
Skywriting script that performs all iterations and convergence testing
in a single CIEL job.

In this experiment, we compare the performance of the two versions by
running 5 iterations of clustering on 20 workers. Each task takes 64 MB
of input---80,000 dense vectors, each containing 100 double-precision
values---and k = 100 cluster centres. We increase the number of tasks
from 20 to 100, in multiples of the cluster size. As before, both
systems use identical data formats and execute an identical
computational kernel. Figure 8(a) compares the per-iteration execution
time for the two versions. For each job size, CIEL is faster than
Hadoop, and the difference ranges between 113 and 168 seconds. To
investigate this difference further, we now analyse the task execution
profile.

Figure 8(b) shows the cluster utilisation as a function of time for the
5 iterations of 100 tasks. From this figure, we can compute the average
cluster utilisation: i.e.~the probability that a worker is assigned a
task at any point during the job execution. Across all job sizes, CIEL
achieves \(89 \pm 2\%\) average utilisation, whereas Hadoop achieves
84\% utilisation for 100 tasks (and only 59\% utilisation for 20 tasks).
The Hadoop utilisation drops to 70\% at several points when there is
still runnable work, which is visible as troughs or ``noise'' in the
utilisation time series. This scheduling delay is due to Hadoop's
polling-based implementation of task dispatch.

CIEL also achieves higher utilisation in this experiment because the
task duration is less variable. The execution time of k-means is
dominated by the map phase, which computes k Euclidean distances for
each data point. Figure 8(c) shows the cumulative distribution of map
task durations, across all k-means experiments. The Hadoop distribution
is clearly bimodal, with 64\% of the tasks being ``fast'' (
\(\mu = 130.9\) , \(\sigma = 3.92\) ) and 36\% of the tasks being
``slow'' ( \(\mu = 193.5\) , \(\sigma = 3.71\) ). By contrast, all of
the CIEL tasks are ``fast'' ( \(\mu = 134.1\) , \(\sigma = 5.05\) ). On
closer inspection, the slow Hadoop tasks are non-data-local: i.e.~they
read their input from another HDFS data node. When computing an
iterative job such as k-means, CIEL can use information about previous
iterations to improve the performance of subsequent iterations. For
example, CIEL preferentially schedules tasks on workers that consumed
the same inputs in previous iterations, in order to exploit data that
might still be stored in the page cache. When a task reads its input
from a remote worker, CIEL also updates the object table to record that
another replica of that input now exists. By contrast, each iteration on
Hadoop is an independent job, and Hadoop does not perform cross-job
optimisations, so the scheduler is less able to exploit data locality.

In the CIEL version, a Skywriting task performs a convergence test and,
if necessary, spawns a subsequent iteration of k-means. However,
compared to the data-intensive map phase, its execution time is
insignificant: in the 100-task experiment, less than 2\% of the total
job

\textit{[Image omitted: images/f2dc5308c338bfd174cddeab06c7a30307f35703dbc89480cf69d022ec431c6d.jpg]}

flowchart

\begin{Shaded}
\begin{Highlighting}[]
\NormalTok{graph TD}
\NormalTok{    A[" "] {-}{-}\textgreater{} C[" "]}
\NormalTok{    B[" "] {-}{-}\textgreater{} C[" "]}
\NormalTok{    C {-}{-}\textgreater{} D[" "]}
\NormalTok{    style C fill:\#f9f,stroke:\#333}
\end{Highlighting}
\end{Shaded}

\begin{enumerate}
\def\labelenumi{(\alph{enumi})}
\tightlist
\item
  Smith-Waterman
\end{enumerate}

\textit{[Image omitted: images/71c5c6859a6037163a06d04fd34935a11b0b12a8f8734f48dd3d0f2db879fefd.jpg]}

natural\_image

Simple geometric diagram showing a triangle with an internal square and
a small circle, no text or symbols present.

\begin{enumerate}
\def\labelenumi{(\alph{enumi})}
\setcounter{enumi}{1}
\tightlist
\item
  Binomial options pricing
\end{enumerate}

Figure 9: Smith-Waterman (§6.3) and BOPM (§6.4) are dynamic programming
algorithms, with macro-level (partition) and micro-level (element)
dependencies.\\
\textit{[Image omitted: images/54b586d125e2abb12b6e2a7ab9217a5c4f4ff797acd1547cbae191535b22d37f.jpg]}

line

{\def\LTcaptype{none} % do not increment counter
\begin{longtable}[]{@{}
  >{\raggedright\arraybackslash}p{(\linewidth - 10\tabcolsep) * \real{0.1391}}
  >{\raggedright\arraybackslash}p{(\linewidth - 10\tabcolsep) * \real{0.1722}}
  >{\raggedright\arraybackslash}p{(\linewidth - 10\tabcolsep) * \real{0.1722}}
  >{\raggedright\arraybackslash}p{(\linewidth - 10\tabcolsep) * \real{0.1722}}
  >{\raggedright\arraybackslash}p{(\linewidth - 10\tabcolsep) * \real{0.1722}}
  >{\raggedright\arraybackslash}p{(\linewidth - 10\tabcolsep) * \real{0.1722}}@{}}
\toprule\noalign{}
\begin{minipage}[b]{\linewidth}\raggedright
Time since start (s)
\end{minipage} & \begin{minipage}[b]{\linewidth}\raggedright
Number of chunks (10 × 10)
\end{minipage} & \begin{minipage}[b]{\linewidth}\raggedright
Number of chunks (20 × 20)
\end{minipage} & \begin{minipage}[b]{\linewidth}\raggedright
Number of chunks (30 × 30)
\end{minipage} & \begin{minipage}[b]{\linewidth}\raggedright
Number of chunks (40 × 40)
\end{minipage} & \begin{minipage}[b]{\linewidth}\raggedright
Number of chunks (50 × 50)
\end{minipage} \\
\midrule\noalign{}
\endhead
\bottomrule\noalign{}
\endlastfoot
0 & \textasciitilde10 & \textasciitilde10 & \textasciitilde10 &
\textasciitilde10 & \textasciitilde10 \\
1654 & \textasciitilde10 & \textasciitilde10 & \textasciitilde10 &
\textasciitilde10 & \textasciitilde10 \\
4122 & \textasciitilde10 & \textasciitilde10 & \textasciitilde10 &
\textasciitilde10 & \textasciitilde10 \\
\end{longtable}
}

Figure 10: Smith-Waterman cluster utilisation against time, for
different block granularities. The best performance is observed with
\(30 \times 30\) blocks.

execution time is spent running Skywriting tasks. The Skywriting
execution time is dominated by communication with the master, as the
script sends a new task descriptor to the master for each task in the
new iteration.

\section{6.3 Smith-Waterman}\label{smith-waterman}

In this experiment, we evaluate strategies for parallelising the
Smith-Waterman sequence alignment algorithm {[}36{]}. For strings of
size m and n, the algorithm computes mn elements of a dynamic
programming matrix. However, since each element depends on three
predecessors, the algorithm is not embarrassingly parallel. We divide
the matrix into blocks---where each block depends on values from its
three neighbours (Figure 9(a))---and process one block per task.

We use CIEL to compute the alignment between two 1 MB strings on 20
workers. Figure 10 shows the cluster utilisation as the block
granularity is varied: a granularity of \(m \times n\) means that the
computation is split into mn blocks. For \(10 \times 10\) (the most
coarse-grained case that we consider), the maximum degree of parallelism
is 10, because the dependency structure limits the maximum achievable
parallelism to the length of the antidiagonal in the block matrix.
Increasing the number of blocks to \(20 \times 20\) allows CIEL to
achieve full utilisation briefly, but performance remains poor because
the majority of the job duration is spent either ramping up to or down
from full utilisation. We observe the best performance for
\(30 \times 30\) , which ramps up to full utilisation more quickly than
coarser-grained configurations, and maintains full utilisation for an
extended period, because there are more runnable tasks than workers.
Increasing the granularity beyond \(30 \times 30\) leads to poorer
overall performance, because the overhead of task dispatch becomes a
significant fraction of task duration. Furthermore, the scheduler cannot
dispatch tasks quickly enough to maintain full utilisation, which
appears as ``noise'' in Figure 10.

\textit{[Image omitted: images/71201f9cd34f8e8238553801e4a882b3f6d5dd770808f0f0563022f0948d7220.jpg]}

line

{\def\LTcaptype{none} % do not increment counter
\begin{longtable}[]{@{}lllll@{}}
\toprule\noalign{}
Number of tasks & 1600k & 800k & 400k & 200k \\
\midrule\noalign{}
\endhead
\bottomrule\noalign{}
\endlastfoot
0 & 0 & 0 & 0 & 0 \\
20 & 10 & 8 & 7 & 5 \\
40 & 15 & 12 & 9 & 5 \\
60 & 20 & 15 & 9 & 4 \\
80 & 23 & 16 & 8 & 3 \\
100 & 24 & 16 & 8 & 3 \\
\end{longtable}
}

Figure 11: Speedup of BOPM (§6.4) on 47 workers as the number of tasks
is varied and the resolution is increased.

\section{6.4 Binomial options pricing}\label{binomial-options-pricing}

We now consider another dynamic programming algorithm: the binomial
options pricing model (BOPM) {[}17{]}. BOPM computes a binomial tree,
which can be represented as an upper-triangular matrix, P. The rightmost
column of P can be computed directly from the input parameters, after
which element \(p_{i,j}\) depends on \(p_{i,j+1}\) and \(p_{i+1,j+1}\) ,
and the result is the value of \(p_{1,1}\) . We achieve parallelism by
dividing the matrix into row chunks, creating one task per chunk, and
streaming the top row of each chunk into the next task. Figure 9(b)
shows the element-and chunk-level data dependencies for this algorithm.

BOPM is not an embarrassingly parallel algorithm. However, we expect
CIEL to achieve some speedup, since rows of the matrix can be computed
in parallel, and we can use streaming tasks ( \(\S5.3\) ) to obtain
pipelined parallelism. We can also achieve better speedup by increasing
the resolution of the calculation: the problem size

\textit{[Image omitted: images/94464e402fc8d61373c82d49abaaab54526a327b5d627b7a2a9c8ed0b0229e54.jpg]}

line

{\def\LTcaptype{none} % do not increment counter
\begin{longtable}[]{@{}lll@{}}
\toprule\noalign{}
Time since start (s) & No failure & Master failure \\
\midrule\noalign{}
\endhead
\bottomrule\noalign{}
\endlastfoot
0 & 100\% & 100\% \\
185 & 100\% & 100\% \\
215 & 100\% & 100\% \\
\end{longtable}
}

Figure 12: Cluster utilisation for three iterations of an iterative
algorithm ( \(\S6.5\) ). In the lower case, the primary master fails
over to a secondary at the beginning of the second iteration. The total
downtime is 7.7 seconds.

\begin{enumerate}
\def\labelenumi{(\alph{enumi})}
\setcounter{enumi}{13}
\tightlist
\item
  is inversely proportional to the time step ( \(\Delta t\) ), and the
  serial execution time increases as \(O(n^{2})\) .
\end{enumerate}

Figure 11 shows the parallel speedup of BOPM on a 47-worker CIEL
cluster. We vary the number of tasks, and increase n from
\(2 \times 10^{5}\) to \(1.6 \times 10^{6}\) . As expected, the maximum
speedup increases as the problem size grows, because the amount of
independent work in each task grows. For \(n = 2 \times 10^{5}\) the
maximum speedup observed is \(4.9 \times\) , whereas for
\(n = 1.6 \times 10^{6}\) the maximum speedup observed is
\(23.8 \times\) . After reaching the maximum, the speedup decreases as
more tasks are added, because small tasks suffer proportionately more
from constant per-task overhead. Due to our streaming implementation,
the minimum execution time for a stream consumer is approximately one
second. We plan to replace our simple, polling-based streaming
implementation with direct TCP sockets, which will decrease the per-task
overhead and improve the maximum speedup.

\section{6.5 Fault tolerance}\label{fault-tolerance-1}

Finally, we conducted an experiment in which master fail-over was
induced during an iterative computation. Figure 12 contrasts the cluster
utilisation in the non-failure and master-failure cases, where the
master fail-over occurs at the beginning of the second iteration.
Between the failure of the primary master and the resumption of
execution, 7.7 seconds elapse: during this time, the secondary master
must detect primary failure, contact all of the workers, and wait until
the workers register with the secondary. Utilisation during the second
iteration is poorer, because some tasks must be replayed due to the
failure. The overall job execution time increases by 30 seconds, and the
original full utilisation is attained once more in the third iteration.

\section{7 Alternative approaches}\label{alternative-approaches}

CIEL was inspired primarily by the MapReduce and Dryad distributed
execution engines. However, there are several different and
complementary approaches to large-scale distributed computing. In this
section, we briefly survey the related work from different fields.

\section{7.1 High performance computing
(HPC)}\label{high-performance-computing-hpc}

The HPC community has long experience in developing parallel programs.
OpenMP is an API for developing parallel programs on shared-memory
machines, which has recently added support for task parallelism with
dependencies {[}7{]}. In this model, a task is a C or Fortran function
marked with a compiler directive that identifies the formal parameters
as task inputs and outputs. The inputs and outputs are typically large
arrays that fit completely in shared memory. OpenMP is more suitable
than CIEL for jobs that share large amounts of data that is frequently
updated on a fine-grained basis. However, the parallel efficiency of a
shared memory system is limited by interconnect contention and/or
non-uniform memory access, which limits the practical size of an OpenMP
job. Nevertheless, we could potentially use OpenMP to exploit
parallelism within an individual multi-core worker.

Larger HPC programs typically use the Message Passing Interface (MPI)
for parallel computing on distributed memory machines. MPI provides
low-level primitives for sending and receiving messages, collective
communication and synchronisation {[}30{]}. MPI is optimised for
low-latency supercomputer interconnects, which often have a
three-dimensional torus topology {[}35{]}. These interconnects are
optimal for problems that decompose spatially and have local
interactions with neighbouring processors. Since these interconnects are
highly reliable, MPI does not tolerate intermittent message loss, and so
checkpointing is usually used for fault tolerance. For example, Piccolo,
which uses MPI, must restart an entire computation from a checkpoint if
an error occurs {[}34{]}.

\section{7.2 Programming languages}\label{programming-languages}

Various programming paradigms have been proposed to simplify or fully
automate software parallelisation.

Several projects have added parallel language constructs to existing
programming languages. Cilk-NOW is a distributed version of Cilk that
allows developers to spawn a C function on another cluster machine and
sync on its result {[}11{]}. X10 is influenced by Java, and provides
finish and async blocks that allow developers to implement more general
synchronisation patterns {[}15{]}. Both implement strict multithreading,
which restricts synchronisation to between a spawned thread

and its ancestor {[}10{]}. While this does not limit the expressiveness
of these languages, it necessitates additional synchronisation in the
implementation of, for example, MapReduce, where non-ancestor tasks may
synchronise.

Functional programming languages offer the prospect of fully automatic
parallelism {[}8{]}. NESL contains a parallel ``apply to each'' operator
(i.e.~a map() function) that processes the elements of a sequence in
parallel, and the implementation allows nested invocation of this
operator {[}9{]}. Glasgow Distributed Haskell contains mechanisms for
remotely evaluating an expression on a particular host {[}33{]}. Though
theoretically appealing, parallel functional languages have not
demonstrated as great scalability as MapReduce or Dryad, which sacrifice
expressivity for efficiency.

\section{7.3 Declarative programming}\label{declarative-programming}

The relational algebra, which comprises a relatively small set of
operators, can be parallelised in time (pipelining) and space
(partitioning) {[}19{]}. Pig and Hive implement the relational algebra
using a DAG of MapReduce jobs on Hadoop {[}32, 37{]}; DryadLINQ and
SCOPE implement it using a Dryad graph {[}14, 39{]}.

The relational algebra is not universal but can be made more expressive
by adding a least fixed point operator {[}5{]}, and this research
culminated in support for recursive queries in SQL:1999 {[}20{]}.
Recently, Bu et al.~showed how some recursive SQL queries may be
translated to iterative Hadoop jobs {[}12{]}.

Datalog is a declarative query language based on first-order logic
{[}13{]}. Recently, Alvaro et al.~developed a version of Hadoop and the
Hadoop Distributed File System using Overlog (a dialect of Datalog), and
demonstrated that it was almost as efficient as the equivalent Java
code, while using far fewer lines of code {[}6{]}. We are not aware of
any project that has used a fully-recursive logic-programming language
to implement data-intensive programs, though the non-recursive Cascalog
language, which runs on Hadoop, is a step in this direction {[}29{]}.

\section{7.4 Distributed operating
systems}\label{distributed-operating-systems}

Hindman et al.~have developed the Mesos distributed operating system to
support ``diverse cluster computing frameworks'' on a shared cluster
{[}25{]}. Mesos performs fine-grained scheduling and fair sharing of
cluster resources between the frameworks. It is predicated on the idea
that no single framework is suitable for all applications, and hence the
resources must be virtualised to support different frameworks at once.
By contrast, we have designed CIEL with primitives that support any form
of computation (though not always optimally), and allow frameworks to be
virtualised at the language level.

\section{8 Conclusions}\label{conclusions}

We designed CIEL to provide a superset of the features that existing
distributed execution engines provide. With Skywriting, it is possible
to write iterative algorithms in an imperative style and execute them
with transparent fault tolerance and automatic distribution. However,
CIEL can also execute any MapReduce job or Dryad graph, and the support
for iteration allows it to perform Pregel- and Piccolo-style
computations.

Our next step is to integrate CIEL primitives with existing programming
languages. At present, only Skywriting scripts can create new tasks.
This does not limit universality, but it requires developers to rewrite
their driver programs in Skywriting. It can also put pressure on the
Skywriting runtime, because all scheduling-related control-flow
decisions must ultimately pass through interpreted code. The main
benefit of Skywriting is that it masks the complexity of
continuation-passing style behind the dereference operator ( \(§4.2\) ).
We now seek a way to extend this abstraction to mainstream programming
languages.

CIEL scales across hundreds of commodity machines, but other scaling
challenges remain. For example, it is unclear how best to exploit
multiple cores in a single machine, and we currently pass this problem
to the executors, which receive full use of an individual machine. This
gives application developers fine control over how their programs
execute, at the cost of additional complexity. However, it limits
efficiency if tasks are inherently sequential and multiple cores are
available. Furthermore, the I/O saving from colocating a stream producer
and its consumers on a single host may outweigh the cost of CPU
contention. Finding the optimal schedule is a hard problem, and we are
investigating simple annotation schemes and heuristics that improve
performance in the common case. The recent work on cluster operating
systems and scheduling algorithms {[}25, 27{]} offers hope that this
problem will admit an elegant solution.

Further information about CIEL and Skywriting, including the source
code, a language reference and a tutorial, is available from the project
website:

http://www.cl.cam.ac.uk/netos/ciel/

\section{Acknowledgements}\label{acknowledgements}

We wish to thank our past and present colleagues in the Systems Research
Group at the University of Cambridge for many fruitful discussions that
contributed to the evolution of CIEL. We would also like to thank
Byung-Gon Chun, our shepherd, and the anonymous reviewers, whose
comments and suggestions have been invaluable for improving the
presentation of this work.

\section{References}\label{references}

{[}1{]} Amazon EC2. http://aws.amazon.com/ec2/.\\
{[}2{]} Apache Hadoop. http://hadoop.apache.org/.\\
{[}3{]} Apache Mahout. http://mahout.apache.org/.\\
{[}4{]} JSON. http://www.json.org/.\\
{[}5{]} AHO, A. V., AND ULLMAN, J. D. Universality of data retrieval
languages. In Proceedings of POPL (1979).\\
{[}6{]} ALVARO, P., CONDIE, T., CONWAY, N., ELMELEEGY, K., HELLERSTEIN,
J. M., AND SEARS, R. BOOM Analytics: Exploring data-centric, declarative
programming for the cloud. In Proceedings of EuroSys (2010).\\
{[}7{]} AYGUADÉ, E., COPTY, N., DURAN, A., HOEFLINGER, J., LIN, Y.,
MASSAIOLI, F., TERUEL, X., UNNIKRISHNAN, P., AND ZHANG, G. The design of
OpenMP tasks. IEEE Trans. Parallel Distrib. Syst. 20, 3 (2009),
404--418.\\
{[}8{]} BACKUS, J. Can programming be liberated from the von Neumann
style?: A functional style and its algebra of programs. Commun. ACM 21,
8 (1978), 613--641.\\
{[}9{]} BLELLOCH, G. E. Programming parallel algorithms. Commun. ACM 39,
3 (1996), 85--97.\\
{[}10{]} BLUMOFE, R. D., AND LEISERSON, C. E. Scheduling multi-threaded
computations by work stealing. J. ACM 46, 5 (1999), 720--748.\\
{[}11{]} BLUMOFE, R. D., AND LISIECKI, P. A. Adaptive and reliable
parallel computing on networks of workstations. In Proceedings of USENIX
ATC (1997).\\
{[}12{]} BU, Y., HOWE, B., BALAZINSKA, M., AND ERNST, M. D. HaLoop:
Efficient iterative data processing on large clusters. In Proceedings of
VLDB (2010).\\
{[}13{]} CERI, S., GOTTLOB, G., AND TANCA, L. What you always wanted to
know about Datalog (and never dared to ask). IEEE Trans. on Knowl. and
Data Eng. 1, 1 (1989), 146--166.\\
{[}14{]} CHAIKEN, R., JENKINS, B., LARSON, P.-A., RAMSEY, B., SHAKIB,
D., WEAVER, S., AND ZHOU, J. SCOPE: Easy and efficient parallel
processing of massive data sets. In Proceedings of VLDB (2008).\\
{[}15{]} CHARLES, P., GROTHOFF, C., SARASWAT, V., DONAWA, C., KIELSTRA,
A., EBCIOGLU, K., VON PRAUN, C., AND SARKAR, V. X10: An object-oriented
approach to non-uniform cluster computing. In Proceedings of OOPSLA
(2005).\\
{[}16{]} CONDIE, T., CONWAY, N., ALVARO, P., HELLERSTEIN, J. M.,
ELMELEEGY, K., AND SEARS, R. MapReduce Online. In Proceedings of NSDI
(2010).\\
{[}17{]} Cox, J. C., Ross, S. A., AND RUBINSTEIN, M. Option pricing: A
simplified approach. Journal of Financial Economics 7, 3 (1979),
229--263.\\
{[}18{]} DEAN, J., AND GHEMAWAT, S. MapReduce: Simplified data
processing on large clusters. In Proceedings of OSDI (2004).\\
{[}19{]} DEWITT, D., AND GRAY, J. Parallel database systems: The future
of high performance database systems. Commun. ACM 35, 6 (1992),
85--98.\\
{[}20{]} EISENBERG, A., AND MELTON, J. SQL: 1999, formerly known as
SQL3. SIGMOD Rec. 28, 1 (1999), 131--138.\\
{[}21{]} EKANAYAKE, J., PALLICKARA, S., AND FOX, G. MapReduce for data
intensive scientific analyses. In Proceedings of eScience (2008).

{[}22{]} GHEMAWAT, S., GOBIOFF, H., AND LEUNG, S.-T. The Google file
system. In Proceedings of SOSP (2003).\\
{[}23{]} GUNDA, P. K., RAVINDRANATH, L., THEKKATH, C. A., YU, Y., AND
ZHUANG, L. Nectar: Automatic management of data and computation in data
centers. In Proceedings of OSDI (2010).\\
{[}24{]} HALSTEAD, JR., R. H. Multilisp: A language for concurrent
symbolic computation. ACM Trans. Program. Lang. Syst. 7, 4 (1985),
501--538.\\
{[}25{]} HINDMAN, B., KONWINSKI, A., ZAHARIA, M., GHODSI, A., JOSEPH, A.
D., KATZ, R., SHENKER, S., AND STOICA, I. Mesos: A platform for
fine-grained resource sharing in the data center. In Proceedings of NSDI
(2011).\\
{[}26{]} ISARD, M., BUDIU, M., YU, Y., BIRRELL, A., AND FETTERLY, D.
Dryad: Distributed data-parallel programs from sequential building
blocks. In Proceedings of EuroSys (2007).\\
{[}27{]} ISARD, M., PRABHAKARAN, V., CURREY, J., WIEDER, U., TALWAR, K.,
AND GOLDBERG, A. Quincy: Fair scheduling for distributed computing
clusters. In Proceedings of SOSP (2009).\\
{[}28{]} MALEWICZ, G., AUSTERN, M. H., BIK, A. J., DEHNERT, J. C., HORN,
I., LEISER, N., AND CZAJKOWSKI, G. Pregel: A system for large-scale
graph processing. In Proceedings of SIGMOD (2010).\\
{[}29{]} MARZ, N. Introducing Cascalog.
http://nathanmarz.com/blog/introducing-cascalog/.\\
{[}30{]} MESSAGE PASSING INTERFACE FORUM. MPI: A message-passing
interface standard. Tech. Rep.~CS-94-230, University of Tennessee,
1994.\\
{[}31{]} MURRAY, D. G., AND HAND, S. Scripting the cloud with
Skywriting. In Proceedings of HotCloud (2010).\\
{[}32{]} OLSTON, C., REED, B., SRIVASTAVA, U., KUMAR, R., AND TOMKINS,
A. Pig Latin: A not-so-foreign language for data processing. In
Proceedings of SIGMOD (2008).\\
{[}33{]} POINTON, R. F., TRINDER, P. W., AND LOIDL, H.-W. The design and
implementation of Glasgow Distributed Haskell. In Proceedings of IFL
(2001).\\
{[}34{]} POWER, R., AND LI, J. Piccolo: Building fast, distributed
programs with partitioned tables. In Proceedings of OSDI (2010).\\
{[}35{]} SCOTT, S. L., AND THORSON, G. M. The Cray T3E network: adaptive
routing in a high performance 3D torus. In Proceedings of HOT
Interconnects (1996).\\
{[}36{]} SMITH, T., AND WATERMAN, M. Identification of common molecular
subsequences. Journal of molecular biology 147, 1 (1981), 195--197.\\
{[}37{]} THUSOO, A., SARMA, J. S., JAIN, N., SHAO, Z., CHAKKA, P.,
ZHANG, N., ANTONY, S., LIU, H., AND MURTHY, R. Hive: A petabyte scale
data warehouse using Hadoop. In Proceedings of ICDE (2010).\\
{[}38{]} WHITE, T. Hadoop: The Definitive Guide. O'Reilly Media, Inc.,
2009.\\
{[}39{]} YU, Y., ISARD, M., FETTERLY, D., BUDIU, M., ÚLFAR ERLINGSSON,
GUNDA, P. K., AND CURREY, J. DryadLINQ: A system for general-purpose
distributed data-parallel computing using a high-level language. In
Proceedings of OSDI (2008).\\
{[}40{]} ZAHARIA, M., BORTHAKUR, D., SEN SARMA, J., ELMELEEGY, K.,
SHENKER, S., AND STOICA, I. Delay scheduling: A simple technique for
achieving locality and fairness in cluster scheduling. In Proceedings of
EuroSys (2010).

\end{document}
